\clearpage
\begin{explanationbox}
\textbf{固定双曲线和 P，让两条割线成对开合：}

\begin{center}
\begin{tikzpicture}[x=1cm,y=1cm,font=\small]
\draw[gray!45,rounded corners=3pt,fill=white] (0,0) rectangle (13.8,7.8);
\node[anchor=west,font=\small\bfseries] at (.4,7.3) {固定 C 与 P，只改变割线斜率：两条弦仍具有同一方向};
\begin{scope}[shift={(1.2,4.8)},scale=.68]
\draw[cncurve,samples=140,domain=-4.7:1.5,variable=\t] plot ({sqrt(2*(1+\t*\t))},\t);
\draw[cnmirror] (2,-4.7)--(2,1.6);
\coordinate (P) at (2,1); \coordinate (M) at (1.4285714285714286,-0.1428571428571428); \coordinate (N) at (3.7142857142857144,-2.4285714285714284);
\draw[cnsec] (M)--(P)--(N); \draw[cnchord] (M)--(N);
\fill (P) circle(2pt) node[above left] {$P$};
\fill (M) circle(2pt) node[left=4pt] {$M$};
\fill (N) circle(2pt) node[right=4pt] {$N$};
\end{scope}
\node at (3.3,6.65) {$m=2$};
\node[red!78!black] at (3.3,.8) {$MN:\ x+y=9/7$};
\begin{scope}[shift={(7.9,4.8)},scale=.68]
\draw[cncurve,samples=140,domain=-4.7:1.5,variable=\t] plot ({sqrt(2*(1+\t*\t))},\t);
\draw[cnmirror] (2,-4.7)--(2,1.6);
\coordinate (P) at (2,1); \coordinate (M) at (1.5294117647058822,0.4117647058823529); \coordinate (N) at (6.235294117647059,-4.294117647058823);
\draw[cnsec] (M)--(P)--(N); \draw[cnchord] (M)--(N);
\fill (P) circle(2pt) node[above left] {$P$};
\fill (M) circle(2pt) node[left=4pt] {$M$};
\fill (N) circle(2pt) node[right=4pt] {$N$};
\end{scope}
\node at (10.0,6.65) {$m=1.25$};
\node[red!78!black] at (10.0,.8) {$MN:\ x+y=33/17$};

\end{tikzpicture}

\end{center}

两图都取 $C:x^2/2-y^2=1$、$P(2,1)$，横纵单位及显示范围相同。绿色割线斜率分别为 $m,-m$，红弦的两个截距不同，但方程都形如 $x+y=t$，斜率始终是 $-1$。

\medskip
\textbf{固定方向从哪里找？} 把 $P$ 沿 $x$ 轴翻到 $Q(2,-1)$，它的切线是 $x+y=1$。首页橙线给出这个方向；这里两条红弦沿它的方向改变位置。

\medskip
\textbf{反射的是直线，不是交点：} 两条割线关于 $x=2$ 对称，双曲线却不关于 $x=2$ 对称，所以不能把 $M$ 直接反射成 $N$。易错图会画出 $M$ 的真正镜像，区分这两件事。

\medskip
\textbf{省掉什么计算？} 若只求弦斜率、判断平行，直接看 $Q$ 点切线即可。要求弦的具体方程时，先保留位置参数 $t$，再用内心、过点或交点范围等条件确定它。

\medskip
\textbf{使用前核对：} 两割线同过双曲线上一个点 $P$，斜率互为相反数，且各有不同于 $P$ 的第二个有限交点。
\end{explanationbox}
